
\documentclass{article}
\usepackage{colortbl}
\usepackage{makecell}
\usepackage{multirow}
\usepackage{supertabular}

\begin{document}

\newcounter{utterance}

\centering \large Interaction Transcript for game `wordle\_withclue', experiment `medium\_frequency\_words\_clue\_no\_critic', episode 7 with Qwen3.5{-}9B.
\vspace{24pt}

{ \footnotesize  \setcounter{utterance}{1}
\setlength{\tabcolsep}{0pt}
\begin{supertabular}{c@{$\;$}|p{.15\linewidth}@{}p{.15\linewidth}p{.15\linewidth}p{.15\linewidth}p{.15\linewidth}p{.15\linewidth}}
    \# & \multicolumn{2}{c}{Player} && \multicolumn{2}{c}{Game Master} \\
    \hline

    \theutterance \stepcounter{utterance}  
    & & & \multicolumn{4}{p{0.6\linewidth}}{
        \cellcolor[rgb]{0.9,0.9,0.9}{
            \makecell[{{p{\linewidth}}}]{
                \texttt{\tiny{[P1$\langle$GM]}}
                \texttt{You are a language wizard who likes to guess words by using the given rules.} \\
\\ 
\texttt{Welcome to Wordle! You have six attempts to guess the target word, a valid English word of five lowercase letters (a{-}z). Please use the tags "explanation:" and "guess:" to provide a concise explanation for each guess.} \\
\\ 
\texttt{To help you make an informed guess, you will receive a clue for the word, such as} \\
\texttt{clue: snowy white.} \\
\\ 
\texttt{Here is an example guess based on the clue:} \\
\texttt{explanation: In the fairy tail Snow White, the girl is killed because she eats a poisoned apple. And the word apple has 5 letters.} \\
\texttt{guess: apple} \\
\\ 
\texttt{After each guess, your answer will be validated, and you will receive feedback indicating which letters are correct (green), which letters are correct but in the wrong position (yellow), and which letters are incorrect (red). This feedback can be useful in determining which letters to include or exclude in your next guess.} \\
\\ 
\texttt{For example, the feedback for "apple" might be:} \\
\texttt{guess\_feedback: a<yellow> p<yellow> p<green> l<yellow> e<red>} \\
\\ 
\texttt{The explanation should contain details about how both the guess\_feedback and clue work together in tandem to formulate a new guess.} \\
\\ 
\texttt{Do not generate any other text.} \\
\\ 
\texttt{Let's begin with your first guess.} \\
\\ 
\texttt{clue: steinbrenner, to the yankees} \\
            }
        }
    }
     \\ \\

    \theutterance \stepcounter{utterance}  
    & \multicolumn{4}{p{0.6\linewidth}}{
        \cellcolor[rgb]{0.561,0.737,0.561}{
            \makecell[{{p{\linewidth}}}]{
                \texttt{\tiny{[P1$\rangle$GM]}}
                \texttt{explanation: The clue "steinbrenner, to the yankees" refers to George Steinbrenner, the owner of the New York Yankees. A famous phrase associated with him is "You're fired," but that doesn't fit the context of ownership or the team name directly in a 5{-}letter word. However, considering the team name "Yankees," the word "yankee" itself is 6 letters. Let's think of a 5{-}letter word related to the Yankees or Steinbrenner. Perhaps "owner"? No, that's too generic. Wait, maybe the word is "boss"? No, that's 4 letters. Let's reconsider. The clue might be hinting at a specific term. Actually, the most direct 5{-}letter word related to the Yankees is "yankee" but it's 6 letters. Hmm. Let's try a different angle. Maybe the word is "team"? No, that's too simple. Wait, perhaps the word is "base"? No. Let's think of words that start with 'y' or are related to baseball. Actually, the clue might be pointing to "yankee" but since it's 6 letters, maybe it's a different word. Wait, the clue says "steinbrenner, to the yankees". Maybe the word is "owner"? No. Let's try "boss" again. No. Wait, perhaps the word is "team"? No. Let's think of 5{-}letter words related to the Yankees. "Yankee" is 6 letters. Maybe the word is "base"? No. Wait, perhaps the word is "ball"? No. Let's try a different approach. The clue might be hinting at a specific term like "owner" but that's 5 letters. Wait, "owner" is 5 letters. But is it related to Steinbrenner? Yes, he was the owner. But is it the best guess? Maybe. Let's try "owner".} \\
\texttt{guess: owner} \\
            }
        }
    }
    & & \\ \\

    \theutterance \stepcounter{utterance}  
    & & & \multicolumn{2}{p{0.3\linewidth}}{
        \cellcolor[rgb]{0.9,0.9,0.9}{
            \makecell[{{p{\linewidth}}}]{
                \texttt{\tiny{[GM$|$GM]}}
                \texttt{attempts = 1} \\
\texttt{target\_word = owner} \\
\texttt{target\_word\_clue = steinbrenner, to the yankees} \\
\texttt{guess = owner} \\
\texttt{guess\_feedback = o<green> w<green> n<green> e<green> r<green>} \\
            }
        }
    }
    & & \\ \\

    \theutterance \stepcounter{utterance}  
    & & & \multicolumn{2}{p{0.3\linewidth}}{
        \cellcolor[rgb]{0.9,0.9,0.9}{
            \makecell[{{p{\linewidth}}}]{
                \texttt{\tiny{[GM$|$GM]}}
                \texttt{game\_result = WIN} \\
            }
        }
    }
    & & \\ \\

\end{supertabular}
}

\end{document}
